\section{Motivation}
It is the aim of this research to simulate this plume overlap concept over a range of atmospheric conditions and aircraft frequencies, to investigate the sensitivity of the climate response. The simulation results can then be used to test the mitigation potential of plume superposition through the application of aircraft trajectory optimisation and formation flight.

\subsection{Aircraft Emissions}
For aircraft to achieve and maintain steady level flight, a propulsion system is often required, providing the necessary thrust force to counteract aerodynamic drag and generate sufficient lift to overcome the aircraft's weight. In the commercial aviation industry, which is responsible for ~88\% of global aviation fuel usage \cite{Gossling2020}, it is commonplace to use a gas turbine propulsion system operating on kerosene-based jet fuel. The stoichiometric combustion of kerosene with air produces the mechanical energy required to operate the engine and generate thrust, however it also emits a number of additional chemical species into the aircraft wake in the process. The emission species which are of great chemical and radiative importance are carbon dioxide (\ce{CO_2}), carbon monoxide (\ce{CO}), water vapour (\ce{H2O}), nitrogen oxides (\ce{NO_x} = \ce{NO} + \ce{NO_2}), sulphur dioxide (\ce{SO_2}) and particulate matter (\ce{\ce{PM}}). 

\subsection{Measuring Aviation Climate Impact}
The most commonly used metric to compare the magnitudes of climate impact from a range of emission species is radiative forcing (RF). RF is defined as the perturbation to the net energy balance of the Earth-atmosphere system due to natural or anthropogenic factors of climate change, measured in watts per square metre [\ce{W/m^2}].

The RF of aviation \ce{CO_2} emissions is direct and well understood. Due to its thermodynamic and photochemical stability, \ce{CO_2} has a relatively long atmospheric lifetime on the order of 100 to 1,000 years. This means that aviation \ce{CO_2} emitted into the atmosphere simply serves to increase its atmospheric concentration, leading to the eventual distribution over global spatial scales and a steady, uniform increase in the planetary greenhouse effect. 
The non-\ce{CO_2} aircraft emission species however, are much more chemically active and hence have much shorter lifetimes and a more localised impact. They primarily contribute to climate change through two indirect radiative forcing mechanisms. Firstly, chemically active species such as \ce{NO_x} and \ce{CO} react with the surrounding air to modify chemical concentrations of key climate agents such as ozone (\ce{O_3}) and methane (\ce{CH_4}). Secondly, emissions of H2O and \ce{PM} are responsible for the formation of condensation trails (contrails) and may even trigger the generation of additional cirrus clouds, both of which can affect the Earth's radiative balance. With the chemical and meteorological state of the atmosphere being highly variable with respect to both location and time, this implies that aircraft non-\ce{CO_2} emissions have a variable climate impact depending on when and where they are released. The spatio-temporal sensitivity of the climate to non-\ce{CO_2} aircraft emissions becomes particularly important when observing the chemical and physical effects that occur in the aircraft exhaust plume. 

  
\subsection{Aircraft Exhaust Plume Treatment}
Following expulsion into the free atmosphere, exhaust species become entrained in the aircraft wake, forming a plume of elevated chemical concentrations which persist for 2 to 12 hours \cite{EPA1992}, before fully dispersing into ambient air. The build-up of non-\ce{CO_2} emissions in the plume gives rise to a number of nonlinear chemical and microphysical effects, which influence the ensuing atmospheric response by altering the net production rates of radiatively active gases and affecting contrail formation and persistence. It is often the case however, that in large-scale atmospheric models, emissions are instantaneously diluted into the volume of the smallest resolved grid cell, with dimensions according to the model's spatial resolution. The instantaneous dilution (ID) approach neglects the plume-scale processes, thus leading to discrepancies in the calculated climate response such as the overestimation of O3 production, \ce{CH_4} and \ce{CO} destruction, and the increased rate of \ce{NO_x} conversion to nitrogen reservoir species \cite{Petry1998}. Furthermore, it is stated in \cite{Fritz2020} that the formation of ice in aircraft exhaust plumes may result in additional heterogeneous chemical reactions, that are not captured in global atmospheric models. 

\subsection{Plume Superposition in High-Density Airspace}
Air traffic levels correlate with aviation demand, giving rise to a spatial and temporal inhomogeneity in the distribution of aircraft, and hence their emissions, across the globe. In a spatial sense, demand tends to be highest around and between densely-populated regions such as the United States, Europe and Asia, meaning flights often take place over mainland airspace and along the flight corridors connecting such regions. In a temporal sense, demand varies both diurnally (time of day) and seasonally (time of year), favouring daylight hours and the summer months \cite{Olsen2013}.

Airspace regions that are particularly susceptible to high demand may see aircraft fly at minimum separation distances from one another. This close proximity flying leads to the superposition of aircraft exhaust plumes, and the accumulation of emissions entrained within them. \cite{Schlager1997} carries out an in-situ investigation into air traffic emissions signatures in the North Atlantic Flight Corridor (NAFC), finding that the superposition of 2-5 plumes lead to the levels of \ce{NO_x}, \ce{SO_2} and \ce{PM} exceeding background concentrations by factors of 30, 5 and 3, respectively. This highlights the severity with which amassing aircraft emission species can alter the local chemical composition of the atmosphere. Such changes to atmospheric conditions can significantly affect the non-\ce{CO_2} climate response of a flight through further chemical and microphysical processing.

\subsection{Atmospheric Response to Plume Superposition}
In \cite{Cameron2013}, the notion of plume superposition and the extent to which it influences the atmospheric response was explored, by comparing results from an expanding plume model against an atmospheric model which assumes instant dilution. The analysis was carried out firstly on a single aircraft at cruise altitude, then on four aircraft flying in the same direction each separated by 1~hr. Observing the single aircraft's plume 10~hrs after emission, it was found that the instantaneous dilution approach overestimated O3 production by 33\%, \ce{CH_4} destruction by 30\% and \ce{CO} destruction by 32\%. When observing the response due to the four overlapping plumes however, the overestimation of O3 production increases to 77\%, \ce{CH_4} destruction to 68\% and \ce{CO} destruction to 74\%. 

The marked difference in measured response for four overlapping plumes compared to a single plume serves as evidence that emissions saturation in superimposed plumes can significantly alter the chemical response of the atmosphere. However, further work needs to be done to clarify the direction and magnitude of the net change in climate impact resulting from these changes to the chemical composition. Moreover, the sensitivity of the chemical and climate response to plume superposition needs to be tested for a range of aircraft transit frequencies and ambient atmospheric conditions, so as to build up a visual representation over the phase space of the atmosphere.

In addition to chemical changes in the atmosphere, plume overlap can also influence the formation and persistence of aircraft contrails and the subsequent generation of aviation-induced cirrus clouds. \cite{Schumann2015} explored the effect of recurrent water vapour emission on contrail formation and persistence at typical aircraft cruising altitudes. It was discovered that continual contrail generation actually leads to dehydration of the surrounding atmosphere and a diminishing radiative forcing effect of \textasciitilde15\% in the affected areas. This outcome provides further impetus to explore the potential effects of plume superposition on the global warming induced by aviation.

\subsection{Mitigation Potential}
The significance of the plume overlap effect and its potential for influencing the atmospheric response brings forth the potential for controlled implementation of this effect, as a means of aviation climate impact mitigation. A near-term fuel burn reduction measure that is currently popular in the literature is the concept of formation flight, which has the potential to provide a complementary non-\ce{CO_2} climate reduction by exploiting the benefits of plume superposition.
Formation flight involves the flight of two or more aircraft in aerodynamic formation, with the follower aircraft positioned in the smooth updraft of the leader aircraft's wake, reducing required lift and thrust, hence reducing fuel and \ce{CO_2} emissions by 5-8\% per trip \cite{Xu2014}. A fortuitous outcome of formation flight is however, the saturation of emissions in the trailing exhaust plumes, which can lead to the aforementioned nonlinear chemical and microphysical response that reduces ozone and contrail formation. 

The non-\ce{CO_2} climate effects of a twin-aircraft formation flight are observed in \cite{Dahlmann2020}, where a climate model was used to determine the changes in \ce{NO_x}-related ozone production and contrail formation processes. The case study found that, despite a 1-3\% increase in flown distance, \ce{CO_2} was reduced by 6\% and \ce{NO_x} was reduced by 11\%. This resulted in a 5\% reduction in ozone production efficiency due to \ce{NO_x}  saturation and a 48\% contrail reduction due to mutual competition for available water vapour, resulting in a total climate impact reduction of approximately 23\%. While part of the alleviated climate impact can be related to the decrease in total emissions due to wake energy retrieval of the follower aircraft, there is also further emphasis due to the saturation of emissions in the overlapping plumes of both aircraft involved. This demonstrates great potential for reducing both \ce{CO_2}- and non-\ce{CO_2}-induced climate effects from aviation, if such a scheme were to be carried out on a global scale.

Fleet-wide implementation of formation flight would however, pose a number of research obstacles, that must be overcome to maximise its cost-benefit potential, whilst ensuring the safe and orderly operation of aircraft in controlled airspace. Optimising aircraft trajectories to minimise fuel burn requires the consideration of nonlinear aircraft performance, wind and weather, payload, fuel load, and constraints set out by air traffic control \cite{Soler2015}. Formation flight adds another variable to the fuel burn optimisation problem, maximising the time spent flying in formation whilst minimising deviation from the true optimal route. This is a research topic covered extensively in the literature  \cite{Ning2011, Hartjes2019}. Optimising formation flight trajectories to minimise climate impact on the other hand, is a relatively novel concept, which requires deeper understanding of the sensitivity of the climate response to emissions under a range of different atmospheric states. 
